%-------------------------------------------------------------------------------
\section{Introduction}
%-------------------------------------------------------------------------------
Recent breakthroughs in Transformer-based models have driven the widespread deployment of high-performance Large Language Models (LLMs) via Model-as-a-Service (MaaS) platforms, where users receive outputs through remote inference APIs. This architecture inherently obscures the generation process: users cannot verify whether the provider substituted a cheaper model, altered outputs before returning them, or silently changed the model during an update. This problem of model execution integrity is particularly critical in high-stakes domains such as finance, healthcare, and security.

Disclosing model parameters for verification is not a viable remedy: model weights embody substantial investments in data, compute, and engineering, and are protected as proprietary intellectual property. A fundamental dilemma thus arises: users demand integrity guarantees for the model's outputs, while providers must preserve the privacy of their parameters.

Zero-knowledge proofs (ZKPs) provide a principled way to reconcile these conflicting requirements. By leveraging ZKPs, a prover (the model provider) can cryptographically prove that an output $y$ is the result of evaluating the model on the user's input $x$ with a secret weight $W$, i.e., $f(x, W) = y$, without revealing any information about $W$ itself. This enables verifiable inference while maintaining the confidentiality of model parameters.

However, making ZKP-based verifiable AI practically deployable requires efficiently proving computationally intensive deep learning operations, and in modern Transformer architectures the primary bottleneck is self-attention. For an input sequence of length $n$, self-attention forms an $n \times n$ attention matrix and applies a row-wise softmax over its entirety, incurring $\mathcal{O}(n^2 d)$ linear operations plus an additional $\mathcal{O}(n^2)$ nonlinear cost (Section~3.1).

Probabilistic techniques such as Freivalds' algorithm~\cite{Freivalds} and Thaler's randomized checks~\cite{Thaler} reduce matrix-multiplication verification from $\mathcal{O}(n^2 d)$ to $\mathcal{O}(n^2)$, yet the verification cost remains anchored at $\mathcal{O}(n^2)$. This quadratic term becomes increasingly prohibitive as context windows grow: BERT~\cite{BERT} uses 512 tokens and GPT-3~\cite{GPT3} 2,048, while frontier models now reach 400K (GPT-5~\cite{GPT5}) and 10M (Llama~4 Scout~\cite{llama4}). As $n$ grows, quadratic verification cost renders proof generation for realistic sequence lengths impractical.

Meanwhile, attention matrices often exhibit approximate low-rank structure that several efficient-attention architectures exploit to reduce inference cost to $\mathcal{O}(n)$: Nystr\"{o}mformer~\cite{Nystromformer} reconstructs the attention map from $m$ landmark tokens via the Nystr\"{o}m method, Linformer~\cite{Linformer} projects keys and values into a $k$-dimensional space, and SOFT~\cite{SOFT} replaces softmax with a Gaussian kernel and applies a Nystr\"{o}m decomposition. Crucially, this efficiency does not sacrifice predictive quality~\cite{Nystromformer,Linformer}. Yet, to the best of our knowledge, verifiable-AI research has not exploited the low-rank structure of attention to reduce proof generation costs, keeping proof complexity at $\mathcal{O}(n^2)$. We bridge this gap with a general framework that combines low-rank factorization with probabilistic verification to eliminate the quadratic bottleneck in verifying self-attention.

\subsection{Our Contribution}
In this work, we propose \name, a verification framework that eliminates the dominant $\mathcal{O}(n^2)$ bottleneck in proving self-attention by exploiting the low-rank factorization common to efficient attention mechanisms. Our main contributions are summarized as follows:
\begin{itemize}
    \item \textbf{A first verification framework for low-rank self-attention.} To the best of our knowledge, no prior verifiable-inference work exploits the low-rank structure of attention. We formalize a common verification relation for attention mechanisms that factor the $n \times n$ attention matrix through $m \ll n$ intermediate dimensions, reducing the circuit complexity of attention verification to $\mathcal{O}(n)$.
    \item \textbf{Constraint-efficient verification via Freivalds' checks and CC-SNARKs.} All large matrix products are computed outside the circuit and verified inside it through Freivalds' randomized projections, so the circuit performs only vector inner products. A commit-carrying SNARK commits to the weights and intermediate matrices outside the circuit, binding every proof to them at a fraction of the constraint cost.
    \item \krev{\textbf{Implementation and evaluation.} We implement and benchmark \name with three instantiations---Nystr\"{o}mformer~\cite{Nystromformer}, Linformer~\cite{Linformer}, and SOFT~\cite{SOFT}---covering landmark-based, projection-based, and kernel-based factorizations. On the Tiny configuration all three reach the same maximum provable sequence length ($n{=}1024$, a $2{\times}$ extension over standard BERT-Freivalds). A multi-spec sweep of Nystr\"{o}mformer across five BERT-family configurations extends this to $4{\times}$.}
\end{itemize}
